% Folder 57, batch 13 (archivists' pages 241-260).
% Pass: Opus 5.5 (claude-opus-5-5), 2026-10-03 — first pass, unchecked against the pages by a human.
% Skipped: 241, 243 (Bourbaki typescript n° 227, "L'application rationnelle définie par
% un système linéaire"), 246 (blank), 248, 250, 252, 260 (typescript pages of a Bourbaki
% draft on commutative algebra, versos), 254 (four-line private letter draft, no mathematics),
% 256, 258 (typescript "Bourbaki's Diktät", programme of a Bourbaki congress).
% The manuscript runs on the versos of these typescripts: 242, 244, 245, 247, 249, 251,
% 253, 255, 257, 259. Connective prose is fast and largely illegible; formulas carry.
\documentclass[11pt,a4paper]{article}
\input{../preamble/grothendieck}

\folder{57}
\batch{13}
\pages{241}{260}
\foldertitle{Picard : tapuscrit annoté (s.d.), lettres (s.d., 1962).}
\dating{1962-[vers 1968]}
\watermark{Édition de démonstration}

\begin{document}

\page{242}
\note{le raisonnement commence avant ce lot ; la notation ($X_n$, $k'/k$, $M_n$, M.L.) est fixée plus haut}

\[
\operatorname{Pic}(X_n) \xrightarrow{\ \sim\ } \underline{\operatorname{Pic}}_{X_n/k}(k)
\]

et \uncertain{donc} M.L. \ill{}.

\note{le passage qui suit, de « Dans le cas général » à « $X_0$ », est encadré et barré de deux traits obliques}
\struck{Dans le cas général, notons \struck{\ill{}} \uncertain{ceux} [où $k'/k$ est une extension \uncertain{radicielle} de \uncertain{façon} \ill{} \ill{} \ill{} \uncertain{contenue} \ill{} \uncertain{car.} $0$ \uncertain{considérée}, mais on a \ill{} \ill{}],}

\[
\text{\struck{$\operatorname{Pic}(X_m) \longrightarrow \operatorname{Pic}(X_n)$}}
\]

\struck{(\ill{} $m \geq n$) est \uncertain{toujours} \ill{} \ill{}. Le \uncertain{noyau} de \ill{} est \ill{} \ill{} \ill{} \ill{} $m = n+1$, \ill{} \ill{} \ill{} \ill{} \ill{} \ill{} \uncertain{groupe} \ill{}} \struck{$\subset H^2(X_0, F^{(n)})$, $F^{(n)}$} \struck{\ill{} \ill{} \ill{} \ill{} $X_0$.}

\uncertain{Sous} \struck{\ill{}} $X_0 \otimes_k k'$ \uncertain{satisfait} à la condition [\uncertain{énoncée} plus haut] : on a une suite exacte

\marginal{\struck{vérifions ??} À expliciter}

\[
0 \longrightarrow \operatorname{Pic}(X_n) \longrightarrow \underline{\operatorname{Pic}}_{X_n/k}(k) \longrightarrow \check{H}^2(k'/k, M_n)
\]

où $M_n$ est le \uncertain{schéma en} groupes [\uncertain{commutatif} \ill{} \struck{\ill{}} qui \uncertain{décrit} les \uncertain{unités} de la \uncertain{$k$-algèbre} \struck{\ill{}} $H^0(X_n, \mathcal{O}_{X_n})$]. On note que les \uncertain{systèmes} \uncertain{projectifs} \ill{} $(M_n)$ satisfont M.L. [\add{\uncertain{de type fini sur} $k$ !} \uncertain{groupes} \struck{\ill{}} \struck{\ill{}} \uncertain{unipotents} \struck{\ill{}} \uncertain{tous} \uncertain{tous} \ill{} \emph{\uncertain{centraux}} \ill{}]

\page{244}
(\uncertain{$\geq n_0$}) \uncertain{donc} \uncertain{pour} $n$ \uncertain{grand}, et $m \geq n$,

\[
\check{H}^*(k'/k, M_m) \xrightarrow{\ \sim\ } \check{H}^*(k'/k, M_n)
\]

\struck{Or on a un homomorphisme \ill{}} \struck{D'où facilement} \ill{} \ill{} :
\uncertain{Soit} $\xi_n \in \underline{\operatorname{Pic}}_{X_n/k}(k)$ \ill{} \ill{} \ill{} \ill{} \ill{}, \uncertain{et} \ill{} $\xi_m$ \uncertain{dans} $\underline{\operatorname{Pic}}_{X_m/k}(k)$ \uncertain{provenant} de $\operatorname{Pic}(X_m)$, alors il provient \uncertain{même} \ill{} de $\operatorname{Pic}(X_n)$. Ceci, et le fait que $(\underline{\operatorname{Pic}}_{X_n/k}(k))$ satisfait à M.L. (A.R.), montre qu'il en est de même de $(\operatorname{Pic}(X_n))_n$, et le \uncertain{fait} \ill{} \ill{} \ill{} $n$, \struck{\ill{}} \ill{} \struck{\ill{}}

et \ill{} \add{$n'$} \ill{} \uncertain{tel} que
$\operatorname{Im}(\underline{\operatorname{Pic}}_{X_{n'}/k}(k) \to \underline{\operatorname{Pic}}_{X_n/k}(k))$
\uncertain{soit} \uncertain{universel}, \uncertain{alors}
$\operatorname{Im}(\operatorname{Pic}(X_{n'}) \to \operatorname{Pic}(X_n))$
\uncertain{est} \uncertain{universel}. OK.

\note{en bas à gauche, un schéma ; les flèches verticales sont tracées vers le haut, comme ici. Il porte en outre les éléments $\xi_m$, $\xi_{n'}$, $\xi_n$ et $\lambda_m$, reliés par des flèches ondulées ; une flèche $\underline{\operatorname{Pic}}_{X_m/k}(k) \leftarrow \check{H}^2(k'/k, M_m)$ y est barrée}

\begin{tikzcd}[column sep=small, row sep=small]
\operatorname{Pic}(X_m) \arrow[r] & \underline{\operatorname{Pic}}_{X_m/k}(k) \\
\operatorname{Pic}(X_{n'}) \arrow[r] \arrow[u] & \underline{\operatorname{Pic}}_{X_{n'}/k}(k) \arrow[u] \\
\operatorname{Pic}(X_n) \arrow[r] \arrow[u] & \underline{\operatorname{Pic}}_{X_n/k}(k) \arrow[u]
\end{tikzcd}

\emph{Corollaire}. Pour que $\mathfrak{X}/Y$ soit \emph{projectif}, il faut et il suffit que pour tout $n$, $X_n/Y_n$ le soit.

\page{245}
\note{feuillet écrit tête-bêche, barré d'un trait oblique, sans lien apparent avec le reste du lot}
\emph{Question} : Soit $H \subset G$, $H$ \uncertain{simple} \ill{}, \uncertain{contenant} un \struck{\ill{}} conjugué de $T$, \struck{Alors} géométrique. Soit $S$ un \uncertain{tore} dans $H$. Alors $S$ \uncertain{est} \uncertain{contenu} dans un \uncertain{tore} \uncertain{maximal} \ill{} de $H$.

\page{247}
\section{§1 Un théorème de géométrie analytique.}

Rappelons d'abord deux résultats connus.

\emph{Lemme} 1.1. Soient $X$ un espace \uncertain{topologique} \ill{}, $Y$ une partie fermée de $X$ \uncertain{ayant} un syst. fondamental de voisinages paracompacts, $F$ un faisceau abélien sur $X$, $F|Y$ le faisceau induit sur $Y$. Alors l'homomorphisme \uncertain{canonique} \struck{\ill{}} suivant est un isomorphisme :

\[
\varinjlim_{U} H^*(U, F) \longrightarrow H^*(Y, F|Y)
\]

[où dans la limite inductive, $U$ \uncertain{parcourt} l'ensemble des voisinages de $Y$ dans $X$]

\emph{Corollaire} 1.2. Soit $f : X \to Y$ une application continue \add{\uncertain{fermée}} \struck{\ill{}}, avec $Y$ loc. compact. \struck{Soit} $y \in Y$, $X_y = f^{-1}(\{y\})$, $F$ un faisceau abélien sur $X$, $F|X_y$ le faisceau induit sur $X_y$. Alors on a un isomorphisme canonique

\[
R^*f_*(F)_y \simeq H^*(X_y, F|X_y)
\]

\page{249}
Voici un résultat \struck{\ill{}} \uncertain{profond}, fondamental \struck{\ill{} \ill{} \ill{}} en géométrie analytique :

\emph{Lemme} 1.3 (\uncertain{Grauert}). Soit $f : X \to Y$ un morphisme propre d'espaces analytiques sur $\mathbb{C}$, $F$ un faisceau \uncertain{cohérent} sur $X$. Alors
\begin{enumerate}
  \item[(i)] Les $R^if_*(F)$ sont cohérents.
  \item[(ii)] Soient $y \in Y$, $Y_n = \operatorname{Specan}(\mathcal{O}_y/\mathfrak{m}_y^{n+1})$, \struck{\ill{}} $X_n = X \times_Y Y_n$, $F_n = F \otimes_Y Y_n$, alors le système projectif $(H^i(X_n, F_n))_n$ satisfait M.L., et l'homomorphisme canonique
\[
R^if_*(F)_y^{\wedge} \longrightarrow \varprojlim_n H^i(X_n, F_n)
\]
est un isomorphisme.
\end{enumerate}

\emph{Corollaire} 1.4. Supposons que $y$ soit un point isolé du \uncertain{support} de $R^if_*(F)$, \uncertain{i.e.} que $R^if_*(F)_y$ soit artinien. Alors on a un isomorphisme

\[
R^if_*(F)_y \longrightarrow \varprojlim_n H^i(X_n, F_n)
\]

\page{251}
Considérons maintenant, pour tout $n$, la suite exacte canonique

\[
0 \longrightarrow \mathbb{Z}_{X_n} \xrightarrow{\ \alpha\ } \mathcal{O}_{X_n} \xrightarrow{\ \beta\ } \mathcal{O}^*_{X_n} \longrightarrow 0
\]

d'où une \emph{\uncertain{suite exacte de cohomologie}} \ill{} :

\[
\begin{aligned}
\cdots \longrightarrow H^i(X_n, \mathbb{Z}) &\xrightarrow{\ \alpha^i\ } H^i(X_n, \mathcal{O}_{X_n}) \xrightarrow{\ \beta^i\ } H^i(X_n, \mathcal{O}^*_{X_n}) \\
&\xrightarrow{\ \partial^i\ } H^{i+1}(X_n, \mathbb{Z}) \xrightarrow{\ \alpha^{i+1}\ } H^{i+1}(X_n, \mathcal{O}_{X_n})
\end{aligned}
\]

\note{sous $H^i(X_n,\mathbb{Z})$ et sous $H^{i+1}(X_n,\mathbb{Z})$ il écrit $\simeq H^i(X_0,\mathbb{Z})$, resp. $H^i(X_n,\mathbb{Z})$ ; l'indice de $\mathbb{Z}$ est biffé}

\uncertain{Pour} \ill{} \ill{} \uncertain{variable}, \ill{} \uncertain{homomorphismes} \ill{} \uncertain{forment} \struck{\ill{} \ill{}} \uncertain{un système projectif} de \uncertain{suites exactes}. Je dis qu'il \uncertain{reste} \uncertain{exact} \struck{\ill{} \ill{}} \add{\uncertain{au sens} de (EGA $0_{\mathrm{III}}$)} \uncertain{par passage à la limite}. En effet, \ill{} \ill{} \ill{} \ill{} \ill{}.

\uncertain{Il} \uncertain{suffit} \ill{} les \uncertain{noyaux} \ill{} \uncertain{des} $\alpha^i$, $\beta^i$, $\partial^i$ \uncertain{satisfont} \ill{} \add{\uncertain{et} $\beta^i$} M.L. \uncertain{Pour} \ill{}, \ill{} \ill{} \ill{} \ill{} \ill{} \uncertain{des} \uncertain{systèmes} \uncertain{projectifs} $(H^i(X_n, \mathbb{Z}))_n$ et $(H^i(X_n, \mathcal{O}_{X_n}))_n$ \uncertain{satisfont} [le premier \uncertain{étant} \uncertain{constant}, le \uncertain{deuxième} \uncertain{en vertu} de 1.3 (ii)]. \uncertain{Pour} \uncertain{les} \ill{}, \ill{} \ill{} \uncertain{que} \ill{} \ill{} \ill{} \uncertain{par} les $\operatorname{Ker}(\alpha^{i+1})$, \ill{} \ill{} \ill{} \ill{} \uncertain{derniers} \struck{\ill{}} \uncertain{étant} \uncertain{supp.} \uncertain{projectifs} \ill{}

\page{253}
\note{la phrase commence à la page précédente}
\uncertain{constants}. En effet, (\ill{} ?) $H^i(X_n, \mathbb{Z}) = H^i(X_0, \mathbb{Z})$ est un groupe de type fini, et \uncertain{son} \ill{} \uncertain{par} $\operatorname{Ker} \alpha^i$ \uncertain{est} \ill{}, \uncertain{toujours} [\uncertain{car} \ill{} \ill{} un \uncertain{sous-groupe} de $H^i(X, \mathcal{O}_{X_n})$, qui est un \uncertain{vectoriel} sur $\mathbb{C}$], \struck{\ill{}} il suffit donc de \uncertain{noter} que \ill{} \ill{} \ill{} \ill{}, la \uncertain{suite} \ill{} \uncertain{décroissante} \ill{} \ill{} \uncertain{constante}.

On obtient ainsi, compte tenu de 1.3 (ii) et du \uncertain{fait} qu'une \uncertain{lim.} \ill{} d'une \uncertain{suite} M.L. \uncertain{de} \uncertain{syst.} M.L. \uncertain{est} M.L. :

\emph{\struck{Corollaire} \add{Proposition}} 1.5. Sous les conditions de 1.3,
\begin{enumerate}
  \item[(i)] Le syst. projectif $(H^i(X_n, \mathcal{O}^*_{X_n}))_n$ satisfait M.L.
  \item[(ii)] On a une suite exacte
\end{enumerate}

\[
\begin{aligned}
H^i(X_0, \mathbb{Z}) &\xrightarrow{\ \alpha^{(i)}_\infty\ } R^if_*(\mathcal{O}_X)_y^{\wedge} \xrightarrow{\ \beta^{(i)}_\infty\ } \varprojlim_n H^i(X_n, \mathcal{O}^*_{X_n}) \\
&\xrightarrow{\ \partial^{(i)}_\infty\ } H^{i+1}(X_0, \mathbb{Z}) \xrightarrow{\ \alpha^{(i+1)}_\infty\ } R^{i+1}f_*(\mathcal{O}_X)_y^{\wedge}
\end{aligned}
\]

N.B. \ill{} \ill{} \ill{} \uncertain{une} \ill{} \uncertain{exacte} \ill{} \ill{} \ill{} \ill{} \uncertain{suite} exacte

\[
0 \longrightarrow \mathbb{Z}_X \longrightarrow \mathcal{O}_X \longrightarrow \mathcal{O}^*_X \longrightarrow 0
\]

\ill{} \uncertain{l'homologie} \ill{} de $(R^if_*)_y$, et compte tenu \uncertain{des isomorphismes} de 1.2, \ill{} $R^if_*(\mathbb{Z}_X)_y$ \uncertain{s'identifie} à $H^i(X_0, \mathbb{Z})$. On obtient \uncertain{ainsi} \ill{}

\page{255}
\note{le diagramme (D) suit immédiatement la fin de la page précédente}

(D)

\begin{tikzcd}[column sep=small, row sep=small, nodes={font=\scriptsize}]
H^i(X_0, \mathbb{Z}) \arrow[r] \arrow[d, no head, "="] & R^if_*(\mathcal{O}_X)_y \arrow[r] \arrow[d, "u^i"] & R^if_*(\mathcal{O}^*_X)_y \arrow[r, "\partial^i"] \arrow[d, "w^i"] & H^{i+1}(X_0, \mathbb{Z}) \arrow[r, "\alpha^{i+1}"] \arrow[d, no head, "="] & R^{i+1}f_*(\mathcal{O}_X)_y \arrow[d, "u^{i+1}"] \\
H^i(X_0, \mathbb{Z}) \arrow[r] & R^if_*(\mathcal{O}_X)_y^{\wedge} \arrow[r] & \varprojlim_n H^i(X_n, \mathcal{O}^*_{X_n}) \arrow[r, "\hat{\partial}^i"] & H^{i+1}(X_0, \mathbb{Z}) \arrow[r, "\hat{\alpha}^{i+1}"] & R^{i+1}f_*(\mathcal{O}_X)_y^{\wedge}
\end{tikzcd}

\note{sous ce diagramme, une troisième ligne $H^i(X_0,\mathbb{Z}) \to H^i(X_0,\mathcal{O}^*_{X_0}) \to \cdots$ est barrée d'un zigzag ; la flèche $R^if_*(\mathcal{O}_X)_y \to R^if_*(\mathcal{O}_X)_y^{\wedge}$ est d'abord marquée $\alpha^i$, puis la colonne reçoit les flèches verticales}

\uncertain{commutatif}, \uncertain{qui} \uncertain{montre}, \ill{} \ill{} \uncertain{lemme} \uncertain{des} \uncertain{cinq} :
\begin{enumerate}
  \item[a)] $R^if_*(\mathcal{O}^*_X)_y \to \varprojlim_n H^i(X_n, \mathcal{O}^*_{X_n})$ \uncertain{est} \emph{injectif}, et \uncertain{en} \uncertain{donne} \ill{} \ill{} \ill{} \uncertain{des} $H^i(X_n, \mathcal{O}^*_{X_n})$.
  \item[b)] Pour \struck{\ill{}} que $R^if_*(\mathcal{O}^*_X)_y \to \varprojlim_n H^i(X_n,\mathcal{O}^*_{X_n})$ \uncertain{soit} un \uncertain{isomorphisme}, \ill{} \ill{} \ill{} que $y$ \ill{} \uncertain{isolé} \uncertain{dans} $\operatorname{Supp} R^if_*(\mathcal{O}_X)$.
\end{enumerate}

\marginal{note écrite en travers de la marge gauche, en partie illisible : « \uncertain{Contrôler} \ill{} \ill{} \ill{} \ill{} $R^2f_*(\mathcal{O}^*_X)$ \ill{} \ill{} $H^2(X_0,\mathbb{Z})$ \ill{} \ill{} \ill{} »}

\emph{Corollaire} 1.6. \ill{} \ill{} \ill{} \uncertain{intrinsèque} \ill{} \uncertain{i.e.} \ill{} \ill{} \ill{} \ill{}

\[
\varinjlim_{U \ni y} H^i(f^{-1}(U), \mathcal{O}^*) = R^if_*(\mathcal{O}^*_X)_y \longrightarrow \varprojlim_n H^i(X_n, \mathcal{O}^*_{X_n})
\]

\uncertain{soit} \ill{} \emph{\uncertain{isomorphisme}}.

c) \struck{\ill{} \ill{}} : la \ill{} \ill{} \ill{} \uncertain{égale} \ill{} \ill{}

\[
R^if_*(\mathcal{O}^*_X)_y \longrightarrow \varprojlim_n H^i(X_n, \mathcal{O}^*_{X_n})
\]

\ill{} \uncertain{bijective}.

\note{suit un passage encadré, barré de plusieurs traits ondulés, avec une note marginale elle aussi barrée ; on y lit des bribes : « \uncertain{Si} $\xi \in \varprojlim_n H^i(X_n,\mathcal{O}^*_{X_n})$, il existe un \ill{} », « $H^i(X,\mathcal{O}^*_X)$ », « M.L. (A.R.) », « $\operatorname{Ker} \alpha^{i+1} = \operatorname{Im} \hat{\partial}^i$ » ; en marge, « $\in R^if_*(\mathcal{O}_X)$ », « \ill{} \ill{} dans $H^{i+1}(X_0,\mathbb{Z})$ », « $\operatorname{Im}$ », « $\operatorname{Im}$ », « O.K. »}

\ill{} \uncertain{Vu} \uncertain{injectivité} \ill{} \ill{} \ill{} \ill{} $v = 1$, \uncertain{qui} \uncertain{donne} le \ill{}

\emph{Remarque} 1.7. Soit $f : X \to Y$ une application \uncertain{propre} d'espaces analytiques sur $\mathbb{C}$, $y \in Y$, \uncertain{posons} \uncertain{toujours}
$Y_n = \operatorname{Specan}(\mathcal{O}_y/\mathfrak{m}_y^{n+1})$, $X_n = X \times_Y Y_n$, \struck{\ill{}} \uncertain{donc}
$[\operatorname{Pic}(X_y) = \varinjlim_{U \ni y} \operatorname{Pic}(f^{-1}(U))]$
$= R^1f_*(\mathcal{O}^*_X)_y$,
\struck{\uncertain{Notons} \ill{} l'\ill{}}

\[
\operatorname{Pic}(\hat{X}_y) = \varprojlim_n \operatorname{Pic}(X_n).
\]

\page{257}
\marginal{en travers de la marge gauche : « \ill{} \ill{} \uncertain{plus} \uncertain{précis}, $\exists k$ \uncertain{tel} \uncertain{que} $\operatorname{Im}(H^1(X_{n+k}, \mathcal{O}^*_{X_{n+k}}) \to H^1(X_n,\mathcal{O}^*_{X_n}))$ et $\operatorname{Im}(R^1f_*(\mathcal{O}^*_X)_y \to H^1(X_n, \mathcal{O}^*_{X_n}))$ \uncertain{ont} \uncertain{même} \uncertain{image} \uncertain{dans} \ill{} »}

\note{la nature de cet énoncé (proposition, corollaire) ne se lit pas ; le titre manque}
Considérons l'homomorphisme canonique

\[
\varphi : \operatorname{Pic}(X_y) \longrightarrow \operatorname{Pic}(\hat{X}_y)
\]

\begin{enumerate}
  \item[(i)] Le syst. \uncertain{projectif} $(\operatorname{Pic}(X_n))_{n \in \mathbb{Z}}$ \uncertain{satisfait} M.L. \add{\uncertain{\ill{}} $\operatorname{Pic}(X_n)$, $\rho$}
  \item[(ii)] \ill{} $\operatorname{Pic}(X_y)$ et $\operatorname{Pic}(\hat{X}_y)$ \ill{} $\varphi$ \uncertain{est} \ill{}
  \item[(iii)] \uncertain{Pour} \ill{} que $\varphi$ \uncertain{soit} \ill{} \ill{}, il \ill{} \ill{} \ill{} que $y$ soit \ill{} \uncertain{isolé} \uncertain{du} supp $R^1f_*(\mathcal{O}_X)$ \struck{\ill{}}.
\end{enumerate}

(Cette condition \uncertain{est} \uncertain{vérifiée} \uncertain{en particulier} si $X|(Y - \{y\}) \to Y - \{y\}$ \ill{} \ill{} \uncertain{isomorphisme}).

\note{le passage suivant est entouré et barré de grands traits obliques}
\struck{\emph{Corollaire} 1.10 \ill{} \ill{} \uncertain{analytiques} : les $X_n$ \uncertain{soient} « \uncertain{algébriques} », \ill{} \uncertain{provenant} \ill{} \ill{} \uncertain{propre} $\mathfrak{X}_n$ \uncertain{sur} $Y_n$ \ill{}.} \struck{Soit $\mathfrak{X}$ le \uncertain{schéma} \uncertain{formel} \uncertain{propre} \uncertain{sur}} \struck{$\hat{\mathcal{O}}_y$} \struck{\uncertain{défini} \uncertain{par} \uncertain{les} $\mathfrak{X}_n$. \ill{} \ill{} \ill{} \ill{} \uncertain{Pour} \uncertain{que} $\mathfrak{X}$ \ill{} \uncertain{soit} \ill{}, \uncertain{il} \uncertain{faut} et \ill{} \ill{}} \struck{\uncertain{Conditions} \uncertain{équivalentes} :}

\struck{(i) Il \uncertain{existe} un \uncertain{voisinage} $U$ de $y$ \ill{} \ill{} \ill{} \uncertain{tel} \uncertain{que} $X|U$ \uncertain{soit} \uncertain{projectif} \uncertain{sur} $U$}

\struck{(\ill{} \uncertain{Grauert} \ill{} 1960/61).}

\struck{(ii) \ill{}} \ill{} \ill{} \struck{\ill{}} \note{la suite de ce passage est hachurée et illisible}

\page{259}
\emph{Corollaire} 1.8. Soit $f : X \to Y$ un morphisme propre d'espaces analytiques sur $\mathbb{C}$ \add{$y \in Y$}. Conditions \uncertain{équivalentes} :
\begin{enumerate}
  \item[(i)] $\exists$ un voisinage \uncertain{ouvert} $U$ de $y$ tel que $X|U$ soit \emph{projectif} sur $U$ (\uncertain{cf.} \ill{} \uncertain{Grauert} \ill{})
  \item[(ii)] \uncertain{Pour} tout $n$, \uncertain{posant} $Y_n = \operatorname{Specan}(\mathcal{O}_y/\mathfrak{m}_y^{n+1})$, $X_n = X \times_{Y} Y_n$, \struck{\uncertain{projectif}} \ill{} \ill{} \ill{} \uncertain{algébrique} \uncertain{projectif}.
\end{enumerate}

\emph{Démonstration}. \uncertain{Résulte} de 1.\ill{} \ill{}

\uncertain{Exemples}

\emph{Corollaire} \struck{1.8} 1.9. \struck{\ill{}} Soit $f : X \to Y$ un morphisme propre d'espaces analytiques \struck{\ill{}} sur $\mathbb{C}$, \struck{\ill{} \ill{} \ill{}} \uncertain{tel} que \ill{} \ill{} \ill{} dim $\leq 1$. Alors $f$ est \uncertain{localement} \uncertain{projectif}.

\uncertain{Résultats} \ill{}

\emph{Lemme} 1.9. Un \uncertain{espace} \ill{} analytique \uncertain{compact} de dim $\leq 1$ est \uncertain{projectif}.
\note{le numéro 1.9 sert deux fois sur la page, au corollaire et au lemme}

\end{document}
